\documentclass[11pt]{article}

\usepackage[a4paper,margin=1.12in]{geometry}
\usepackage{amsmath,amssymb,amsthm,mathtools}
\usepackage{microtype}
\usepackage[hidelinks]{hyperref}
\usepackage[nameinlink,capitalise]{cleveref}

\newtheorem{theorem}{Theorem}[section]
\newtheorem{corollary}[theorem]{Corollary}
\newtheorem{proposition}[theorem]{Proposition}
\newtheorem{lemma}[theorem]{Lemma}
\theoremstyle{remark}
\newtheorem{remark}[theorem]{Remark}

\newcommand{\Hh}{\mathcal H}
\newcommand{\Dom}{\mathcal D}
\newcommand{\Ran}{\operatorname{Ran}}
\newcommand{\Gr}{\operatorname{Gr}}
\newcommand{\spec}{\operatorname{spec}}
\newcommand{\Id}{\mathbb I}
\newcommand{\C}{\mathbb C}
\newcommand{\R}{\mathbb R}

\title{Boundary-Phase Locking for Nested Symmetric Operators\\
and Suzuki's Finite-Aperture Hilbert--P\'olya Limit}
\author{Rejected research draft: RH-conditional application; not Lean-complete}
\date{22 July 2026}

\begin{document}
\maketitle

\begin{abstract}
Suzuki recently constructed, at each finite aperture, a symmetric first-order operator of
deficiency indices $(1,1)$ whose self-adjoint extensions have real, simple characteristic zeros.
Under the Riemann hypothesis he also constructed a global translation generator whose spectrum is
the set of ordinates of the nontrivial zeta zeros, and conjectured that a suitable aperture-dependent
boundary phase makes the finite operators converge to this global Hilbert--P\'olya operator in
strong resolvent sense.

We prove an operator-theoretic phase-locking law for every nested exhaustion of a closed symmetric
operator with stable finite deficiency indices.  Graph exhaustion automatically forces convergence
of the deficiency projections.  Their polar transport canonically carries any global boundary
unitary to the finite deficiency spaces.  The corresponding self-adjoint extensions converge in
generalized strong resolvent sense.  Conversely, convergence occurs if and only if the transported
finite boundary unitaries converge to the global one.  For indices $(1,1)$ this reduces the entire
limit problem to convergence of one effective unit complex phase.

Applied to Suzuki's Weil-form Hilbert spaces, the theorem proves under the Riemann hypothesis the
existence of the conjectured phase schedule and classifies all schedules with the same limit.  Every
ordinate of a nontrivial zeta zero is consequently approximated by real zeros of Suzuki's
finite-aperture characteristic functions.  This does not prove the Riemann hypothesis: positivity
of the global Weil Hilbert space and existence of the target self-adjoint generator are precisely
the inputs supplied by that hypothesis.
\end{abstract}

\section{Introduction}

Suzuki's 2026 finite-aperture program recasts the Weil quadratic form in terms of a continuous screw
kernel and a family of self-adjoint differential realizations \cite{Suzuki2026}.  At aperture
$a>0$, the construction produces a Hilbert space $\Hh_a$, a closed symmetric realization $S_a$ of
$i\,d/dx$ with deficiency indices $(1,1)$, and a circle of self-adjoint extensions
\[
 D_{a,\theta},\qquad \theta\in[0,2\pi).
\]
Their spectra are discrete and simple.  Suzuki gives an entire characteristic function
$W(a,\theta;z)$ whose zeros are exactly the eigenvalues of $D_{a,\theta}$, and proves that all of
these zeros are real.

Under the Riemann hypothesis, the global Weil form is positive and its completion $\Hh_\infty$
carries a self-adjoint translation generator $D$.  Through Suzuki's de Branges model, the spectrum
of $D$ is the set $\Gamma$ of ordinates of the nontrivial zeros of the Riemann zeta-function.  The
finite spaces embed into the global space.  Suzuki then states the central expectation that a
boundary schedule $\theta=\theta(a)$ can be selected so that
\begin{equation}\label{eq:Suzuki-expectation}
  D_{a,\theta(a)}\longrightarrow D
  \qquad (a\to\infty)
\end{equation}
in strong resolvent sense \cite[Sections 7.5--7.6]{Suzuki2026}.

Two issues are unresolved in \eqref{eq:Suzuki-expectation}.  First, the operators act in different
Hilbert spaces, so ordinary strong resolvent convergence is not literally defined.  Second, the
finite extension phase must select the particular global extension generated by translation.  The
reality of every finite spectrum does not determine this phase.

We solve both issues.  Our main result is a universal theorem for nested closed symmetric
operators.  If the finite graphs exhaust the global graph and the deficiency indices remain equal
and finite, then the finite deficiency projections converge automatically.  A canonical polar
transport identifies the limiting and finite deficiency spaces.  Transporting the boundary
unitary of a chosen global self-adjoint extension produces finite extensions converging to it in
the generalized strong resolvent sense appropriate to varying subspaces.

The converse is equally important.  Every convergent extension schedule must converge after the
same canonical boundary transport.  Thus the construction does not merely exhibit one phase; it
classifies all successful phases.  At deficiency indices $(1,1)$, if $\zeta_a$ and $\zeta$ are the
finite and global effective phases, then
\begin{equation}\label{eq:scalar-lock-intro}
 D_{a,\theta(a)}\longrightarrow D
 \quad\Longleftrightarrow\quad
 \zeta_a\longrightarrow\zeta.
\end{equation}

There is substantial adjacent literature.  Bjerg characterized convergence of self-adjoint
extensions in terms of convergence of their von Neumann boundary maps when all operators act in
one fixed Hilbert space \cite{Bjerg2024}.  Teschl, Wang, Xie, and Zhou recently streamlined
generalized resolvent convergence for operators on varying Hilbert spaces and established its
functional-calculus and spectral consequences \cite{TeschlEtAl2026}.  Our new step is the automatic
boundary-space convergence and canonical phase selection forced by a nested graph exhaustion,
together with its application to Suzuki's aperture chain.

The universal theorem is developed in \cref{sec:setup,sec:deficiency,sec:phase}.  Spectral
approximation is recorded in \cref{sec:spectral}.  In \cref{sec:Suzuki} we verify the hypotheses
for Suzuki's exact spaces and prove \eqref{eq:Suzuki-expectation} in its rigorous varying-space
form.  The theorem is conditional on RH only at this final specialization; the operator theorem
itself is unconditional.

\section{Nested symmetric exhaustions}\label{sec:setup}

Let $\Hh$ be a complex Hilbert space.  Let $(\Hh_a)_{a>0}$ be an increasing family of closed
subspaces,
\[
 \Hh_a\subset\Hh_b\subset\Hh\qquad (a<b),
\]
and let $P_a:\Hh\to\Hh_a$ be the orthogonal projection.  We assume
\begin{equation}\label{eq:Pa-strong}
 P_a\xrightarrow[a\to\infty]{\mathrm{s}}\Id.
\end{equation}
The inclusion is denoted by $\iota_a:\Hh_a\hookrightarrow\Hh$, so
$\iota_a^*=P_a$.

For each $a$, let $S_a$ be a closed densely defined symmetric operator on $\Hh_a$.  Let $S$ be a
closed densely defined symmetric operator on $\Hh$.  We impose the nested-extension condition
\begin{equation}\label{eq:nested-S}
 S_a\subset S_b\subset S\qquad(a<b),
\end{equation}
where the inclusions use the canonical embeddings of the Hilbert spaces, and the graph-exhaustion
condition
\begin{equation}\label{eq:graph-exhaust}
 \overline{\bigcup_{a>0}\Gr(S_a)}=\Gr(S)
 \quad\text{in }\Hh\oplus\Hh.
\end{equation}
Finally, all deficiency indices are assumed stable:
\begin{equation}\label{eq:stable-deficiency}
 \operatorname{def}(S_a)=\operatorname{def}(S)=(k,k),
 \qquad 1\leq k<\infty.
\end{equation}

For a closed symmetric operator $T$, set
\[
 \mathcal M_+(T)=\Ran(T+i),\qquad
 \mathcal M_-(T)=\Ran(T-i),
\]
and
\[
 \mathcal N_+(T)=\mathcal M_+(T)^\perp=\ker(T^*-i),\qquad
 \mathcal N_-(T)=\mathcal M_-(T)^\perp=\ker(T^*+i).
\]
Orthogonal complements for $S_a$ are taken inside $\Hh_a$.  We abbreviate
\[
 \mathcal M_{a,\pm}=\mathcal M_\pm(S_a),\quad
 \mathcal N_{a,\pm}=\mathcal N_\pm(S_a),\quad
 \mathcal M_\pm=\mathcal M_\pm(S),\quad
 \mathcal N_\pm=\mathcal N_\pm(S).
\]
Let $Q_{a,\pm}$ and $E_{a,\pm}$ denote the orthogonal projections, embedded in $\Hh$, onto
$\mathcal M_{a,\pm}$ and $\mathcal N_{a,\pm}$ respectively.  Let $Q_\pm,E_\pm$ be their global
counterparts.

\begin{remark}[Generalized strong resolvent convergence]\label{rem:GSRC}
If $A_a$ and $A$ are self-adjoint on $\Hh_a$ and $\Hh$, respectively, the convergence used below
is
\begin{equation}\label{eq:GSRC-def}
 \iota_a(A_a+i)^{-1}P_a
 \xrightarrow[a\to\infty]{\mathrm{s}}
 (A+i)^{-1}.
\end{equation}
This is generalized strong resolvent convergence for the projection/inclusion pair.  It is the
correct replacement for ordinary strong resolvent convergence in
\eqref{eq:Suzuki-expectation}; see \cite{TeschlEtAl2026}.
\end{remark}

\section{Automatic convergence of deficiency spaces}\label{sec:deficiency}

The first theorem is the mechanism missing from a phase-selection argument.  It shows that no
separate convergence assumption on deficiency vectors is needed in a stable nested exhaustion.

\begin{lemma}[Closed shifted ranges]\label{lem:closed-range}
If $T$ is closed and symmetric, then $\Ran(T+i)$ and $\Ran(T-i)$ are closed.
\end{lemma}

\begin{proof}
For $x\in\Dom(T)$, symmetry gives
\begin{equation}\label{eq:shifted-graph-norm}
 \|(T\pm i)x\|^2=\|Tx\|^2+\|x\|^2.
\end{equation}
Thus $T\pm i$ is bounded below in the graph norm.  A Cauchy sequence in its range lifts to a graph
Cauchy sequence in $\Dom(T)$, which converges because $T$ is closed.
\end{proof}

\begin{theorem}[Automatic deficiency-projection convergence]\label{thm:deficiency-convergence}
Under \eqref{eq:Pa-strong}--\eqref{eq:stable-deficiency},
\begin{equation}\label{eq:range-projection-convergence}
 Q_{a,\pm}\xrightarrow[a\to\infty]{\mathrm{s}}Q_\pm,
\end{equation}
and
\begin{equation}\label{eq:deficiency-projection-convergence}
 E_{a,\pm}\xrightarrow[a\to\infty]{\|\cdot\|}E_\pm.
\end{equation}

For all sufficiently large $a$, the polar decompositions of
\[
 E_{a,+}|_{\mathcal N_+}:\mathcal N_+\to\mathcal N_{a,+},
 \qquad
 E_{a,-}|_{\mathcal N_-}:\mathcal N_-\to\mathcal N_{a,-}
\]
have unitary polar factors
\begin{equation}\label{eq:J-def}
 J_{a,+}:\mathcal N_+\to\mathcal N_{a,+},
 \qquad
 J_{a,-}:\mathcal N_-\to\mathcal N_{a,-},
\end{equation}
and these satisfy
\begin{equation}\label{eq:J-convergence}
 J_{a,\pm}\longrightarrow\operatorname{incl}_{\mathcal N_\pm\hookrightarrow\Hh}
 \quad\text{in operator norm}.
\end{equation}
\end{theorem}

\begin{proof}
By \eqref{eq:nested-S}, the spaces $\mathcal M_{a,+}$ are increasing and contained in
$\mathcal M_+$.  Let $(S+i)x\in\mathcal M_+$.  Graph exhaustion gives
$x_a\in\Dom(S_a)$ such that
\[
 x_a\to x,\qquad S_ax_a\to Sx.
\]
Therefore
\[
 (S_a+i)x_a\longrightarrow(S+i)x.
\]
It follows that
\[
 \overline{\bigcup_{a>0}\mathcal M_{a,+}}=\mathcal M_+.
\]
The spaces are closed by \cref{lem:closed-range}.  Monotone convergence of orthogonal projections
proves the plus case of \eqref{eq:range-projection-convergence}; the minus case is identical.

Since $\mathcal M_{a,\pm}\subset\Hh_a$, one has
\begin{equation}\label{eq:E-difference}
 E_{a,\pm}=P_a-Q_{a,\pm},
 \qquad
 E_\pm=\Id-Q_\pm.
\end{equation}
Equations \eqref{eq:Pa-strong} and \eqref{eq:range-projection-convergence} give strong convergence
of the deficiency projections.  Their ranks and the ranks of their limits all equal the same finite
number $k$.  Strong convergence of orthogonal projections of constant finite rank to a projection
of that rank is norm convergence, proving \eqref{eq:deficiency-projection-convergence}.

On the finite-dimensional space $\mathcal N_+$, the restriction
$E_{a,+}|_{\mathcal N_+}$ therefore converges in norm to the identity inclusion.  It is injective
for all large $a$.  Its range lies in $\mathcal N_{a,+}$, and equality of dimensions makes it
surjective.  Its polar factor $J_{a,+}$ is consequently unitary.  Continuous dependence of the
polar factor near an isometry gives \eqref{eq:J-convergence}.  The minus case is the same.
\end{proof}

\begin{remark}[Why stable deficiency is load-bearing]
If the finite deficiency dimension exceeds the limiting dimension, a deficiency direction may
escape weakly to infinity; constant-rank projections can then converge strongly to a lower-rank
projection without norm convergence.  Condition \eqref{eq:stable-deficiency} excludes exactly this
boundary leakage.
\end{remark}

\section{Boundary-phase locking and its converse}\label{sec:phase}

For a self-adjoint operator $A$, define the Cayley transform
\begin{equation}\label{eq:Cayley}
 C_A=(A-i)(A+i)^{-1}.
\end{equation}
For a closed symmetric $S$, its partial Cayley isometry
\[
 V:\mathcal M_+\longrightarrow\mathcal M_-
\]
is determined by
\begin{equation}\label{eq:partial-Cayley}
 V((S+i)x)=(S-i)x.
\end{equation}
If $A$ is a self-adjoint extension of $S$, then $C_A$ extends $V$ and decomposes as
\begin{equation}\label{eq:Cayley-block}
 C_A=VQ_+ + WE_+,
\end{equation}
where
\begin{equation}\label{eq:global-boundary-block}
 W:\mathcal N_+\longrightarrow\mathcal N_-
\end{equation}
is unitary.  Working with the Cayley boundary block avoids the convention-dependent minus sign
between Cayley and von Neumann extension parameters.

\begin{theorem}[Canonical boundary-phase locking]\label{thm:phase-lock-existence}
Let $A$ be any self-adjoint extension of $S$, with boundary block $W$ in
\eqref{eq:global-boundary-block}.  For all sufficiently large $a$, set
\begin{equation}\label{eq:locked-boundary}
 W_a=J_{a,-}WJ_{a,+}^*:
 \mathcal N_{a,+}\longrightarrow\mathcal N_{a,-}.
\end{equation}
Let $A_a$ be the self-adjoint extension of $S_a$ with Cayley boundary block $W_a$.  Then
\begin{equation}\label{eq:main-GSRC}
 \iota_a(A_a+i)^{-1}P_a
 \xrightarrow[a\to\infty]{\mathrm{s}}
 (A+i)^{-1}.
\end{equation}
\end{theorem}

\begin{proof}
Let $V_a:\mathcal M_{a,+}\to\mathcal M_{a,-}$ be the partial Cayley isometry of $S_a$.  Since
$S_a\subset S$,
\begin{equation}\label{eq:Va-restriction}
 V_a=V|_{\mathcal M_{a,+}}.
\end{equation}
The embedded finite Cayley transform therefore has the block formula
\begin{equation}\label{eq:embedded-Cayley}
 \iota_aC_{A_a}P_a
 =VQ_{a,+}+W_aE_{a,+}.
\end{equation}
By \cref{thm:deficiency-convergence}, the first term converges strongly to $VQ_+$ and the second
to $WE_+$.  Hence
\begin{equation}\label{eq:Cayley-strong}
 \iota_aC_{A_a}P_a\xrightarrow{\mathrm{s}}C_A.
\end{equation}

The Cayley resolvent identity is
\begin{equation}\label{eq:Cayley-resolvent}
 (A+i)^{-1}=\frac1{2i}(\Id-C_A).
\end{equation}
On the finite space it gives
\[
 \iota_a(A_a+i)^{-1}P_a
 =\frac1{2i}\left(P_a-\iota_aC_{A_a}P_a\right).
\]
Now use \eqref{eq:Pa-strong} and \eqref{eq:Cayley-strong}.
\end{proof}

The next result shows that the construction is exhaustive, not merely sufficient.

\begin{theorem}[Boundary-phase classification]\label{thm:phase-classification}
Let $B_a$ be arbitrary self-adjoint extensions of $S_a$, with Cayley boundary blocks
\[
 Z_a:\mathcal N_{a,+}\longrightarrow\mathcal N_{a,-}.
\]
Then
\begin{equation}\label{eq:arbitrary-GSRC}
 \iota_a(B_a+i)^{-1}P_a
 \xrightarrow[a\to\infty]{\mathrm{s}}
 (A+i)^{-1}
\end{equation}
if and only if
\begin{equation}\label{eq:transported-boundary-convergence}
 J_{a,-}^*Z_aJ_{a,+}\longrightarrow W
 \quad\text{in operator norm}.
\end{equation}
\end{theorem}

\begin{proof}
If \eqref{eq:transported-boundary-convergence} holds, the proof of
\cref{thm:phase-lock-existence} applies with $Z_a$ in place of $W_a$.

Conversely, \eqref{eq:arbitrary-GSRC}, \eqref{eq:Pa-strong}, and the Cayley resolvent identity imply
\[
 \iota_aC_{B_a}P_a\xrightarrow{\mathrm{s}}C_A.
\]
The base blocks $VQ_{a,+}$ converge strongly to $VQ_+$ independently of the extension.  Subtracting
them gives
\begin{equation}\label{eq:boundary-block-strong}
 Z_aE_{a,+}\xrightarrow{\mathrm{s}}WE_+,
\end{equation}
where all finite maps are understood after inclusion into $\Hh$.

Fix $x\in\mathcal N_+$.  Since $J_{a,+}x\to x$ and the boundary blocks are contractions,
\eqref{eq:boundary-block-strong} yields
\[
 Z_aJ_{a,+}x\longrightarrow Wx.
\]
Applying $J_{a,-}^*$ and using \eqref{eq:J-convergence} gives pointwise convergence in the
finite-dimensional space $\mathcal N_+$.  Uniform boundedness and finite dimension upgrade this to
operator-norm convergence, proving \eqref{eq:transported-boundary-convergence}.
\end{proof}

\begin{corollary}[Scalar phase-locking law]\label{cor:scalar-phase}
Suppose $k=1$.  Choose unit vectors $e_\pm\in\mathcal N_\pm$ and write
\[
 We_+=\zeta e_-,\qquad |\zeta|=1.
\]
Every finite boundary block has a unique phase $\zeta_a\in\mathbb T$ defined by
\[
 Z_aJ_{a,+}e_+=\zeta_aJ_{a,-}e_-.
\]
Then \eqref{eq:arbitrary-GSRC} holds if and only if
\begin{equation}\label{eq:scalar-phase-lock}
 \zeta_a\longrightarrow\zeta.
\end{equation}
\end{corollary}

\begin{proof}
Between one-dimensional Hilbert spaces a unitary is multiplication by one unit scalar.
\Cref{thm:phase-classification} becomes \eqref{eq:scalar-phase-lock}.
\end{proof}

\section{Spectral consequences}\label{sec:spectral}

Generalized strong resolvent convergence gives target spectral inclusion, although it does not in
general exclude spectral pollution.

\begin{proposition}[Approximation of isolated target spectrum]\label{prop:spectral-approx}
Assume \eqref{eq:main-GSRC}.  Let $\lambda$ be an isolated point of $\spec(A)$.  If $U$ is any
open interval containing $\lambda$ whose endpoints lie in the resolvent set of $A$, then
\[
 U\cap\spec(A_a)\ne\varnothing
\]
for all sufficiently large $a$.  Equivalently,
\begin{equation}\label{eq:spectral-liminf}
 \spec(A)\subset\liminf_{a\to\infty}\spec(A_a).
\end{equation}
\end{proposition}

\begin{proof}
Choose a nonnegative bounded continuous function $f$ supported in $U$ and satisfying
$f(\lambda)=1$.  Generalized strong resolvent convergence implies
\[
 \iota_af(A_a)P_a\xrightarrow{\mathrm{s}}f(A);
\]
this is the standard continuous-functional-calculus consequence
\cite{TeschlEtAl2026}.  Since $f(A)\ne0$, one has $f(A_a)\ne0$ for all sufficiently large $a$.
Thus $\spec(A_a)$ meets the support of $f$, which is contained in $U$.
\end{proof}

\begin{remark}
Strong resolvent convergence proves the lower spectral limit
\eqref{eq:spectral-liminf}.  It does not prove norm resolvent convergence, equality of spectral
limits, or absence of extra finite-aperture eigenvalues away from $\spec(A)$.
\end{remark}

\section{Suzuki's aperture chain}\label{sec:Suzuki}

We now assume RH and verify the universal hypotheses using Suzuki's constructions.  No new
estimate on primes or zero ordinates is needed.

\subsection{The Hilbert-space exhaustion}

Let $\Hh_a=\Hh(A_a)$ be Suzuki's completion of $C_c^\infty(-a,a)$ in the restricted Weil-form norm,
and let $\Hh_\infty=\Hh(A_\infty)$ be the global Weil Hilbert space.  Under RH the global Weil form
is positive; Suzuki then has $A_a>0$ for every $a$ and takes the finite shift parameter
$\lambda=0$.  On a function supported in $(-a,a)$, the finite and global norms are consequently
both the same Weil quadratic form.  Extension by zero therefore gives isometric embeddings
\begin{equation}\label{eq:Suzuki-space-chain}
 \Hh_a\subset\Hh_b\subset\Hh_\infty
 \qquad(a<b).
\end{equation}
Their images are closed because the finite spaces are complete.  Their union contains
$C_c^\infty(\R)$, which is dense in $\Hh_\infty$ by construction.  Consequently their orthogonal
projections satisfy \eqref{eq:Pa-strong}.

\subsection{The minimal derivative exhaustion}

On $\Hh_a$, Suzuki's minimal derivative is defined on $C_c^\infty(-a,a)$ by
\[
 D_a^0f=i f'.
\]
Let $S_a=\overline{D_a^0}$.  Globally, let
\[
 D_\infty^0f=i f',
 \qquad \Dom(D_\infty^0)=C_c^\infty(\R)\subset\Hh_\infty,
\]
and put $S=\overline{D_\infty^0}$.  Equality of the Weil norms on compact support gives
\[
 S_a\subset S_b\subset S.
\]
Every global test function is supported in a finite interval.  Hence the union of the finite test
graphs is the global test graph, whose closure is $\Gr(S)$.  This proves
\eqref{eq:graph-exhaust}.

Suzuki proves that every finite minimal derivative has deficiency indices $(1,1)$.  Under RH his
global de Branges model identifies $S$ with a symmetric multiplication operator of the same
indices.  Thus \eqref{eq:stable-deficiency} holds with $k=1$.

\subsection{The global translation extension}

Suzuki's translation group on $\Hh_\infty$ has a self-adjoint infinitesimal generator $D$.
It is the extension denoted $D_{\pi/2}$ in the de Branges model, and
\begin{equation}\label{eq:global-spectrum}
 \spec(D)=\Gamma,
\end{equation}
where $\Gamma$ is the discrete set of ordinates of the nontrivial zeta zeros
\cite[Sections 7.2--7.5]{Suzuki2026}.  Let $W:\mathcal N_+\to\mathcal N_-$ be the Cayley boundary
block of this extension.

For all sufficiently large $a$, \cref{thm:phase-lock-existence} defines
\[
 W_a=J_{a,-}WJ_{a,+}^*.
\]
Suzuki parametrizes all finite self-adjoint extensions by $\theta\in[0,2\pi)$, so there is a phase
$\theta(a)$ whose extension has boundary block $W_a$.

\begin{theorem}[Suzuki finite-aperture spectral limit under RH]\label{thm:Suzuki-limit}
Assume RH.  There exists a boundary schedule $\theta(a)$ such that
\begin{equation}\label{eq:Suzuki-GSRC}
 \iota_a(D_{a,\theta(a)}+i)^{-1}P_a
 \xrightarrow[a\to\infty]{\mathrm{s}}
 (D+i)^{-1}.
\end{equation}
Every other phase schedule has the same limit if and only if its effective Cayley phase, transported
by $J_{a,\pm}$, converges to the boundary phase of the global translation generator.
\end{theorem}

\begin{proof}
All hypotheses of \cref{thm:phase-lock-existence,thm:phase-classification} were verified above.
Apply those theorems with $A=D$ and $A_a=D_{a,\theta(a)}$.
\end{proof}

This proves Suzuki's expectation in a precise common-space formulation.  The phase schedule is
canonical after choosing the global extension: it is obtained by orthogonally projecting the global
deficiency lines into the finite spaces and taking their polar normalization.  No fitted spectral
data or target-defined error term enters the construction.

\subsection{Characteristic-zero approximation}\leavevmode\par

Suzuki proves that the zeros of $W(a,\theta;z)$ are exactly the eigenvalues of $D_{a,\theta}$ and
that all are real.  Combining this with \cref{prop:spectral-approx} and
\eqref{eq:global-spectrum} gives the desired spectral approximation.

\begin{corollary}[Every xi ordinate is a finite-aperture spectral limit]\label{cor:Xi-approx}
Assume RH and choose the phase schedule from \cref{thm:Suzuki-limit}.  For every
$\gamma\in\Gamma$ and every $\varepsilon>0$, there exists $a_0$ such that for every $a\geq a_0$
there is a real number $x_a$ satisfying
\begin{equation}\label{eq:characteristic-zero-approx}
 W(a,\theta(a);x_a)=0,
 \qquad
 |x_a-\gamma|<\varepsilon.
\end{equation}
\end{corollary}

\begin{proof}
The global spectrum is discrete, so choose an interval contained in
$(\gamma-\varepsilon,\gamma+\varepsilon)$ with endpoints outside $\Gamma$.  By
\cref{prop:spectral-approx}, this interval meets $\spec(D_{a,\theta(a)})$ for all large $a$.
The characteristic-function theorem identifies such an eigenvalue with a real zero of
$W(a,\theta(a);z)$.
\end{proof}

\section{Relation to Suzuki's characteristic-function proposal}

Suzuki also proposed a locally uniform limit involving
\[
 z^2\frac{\xi(1/2-iz)}{\xi'(1/2-iz)}.
\]
That target is meromorphic, whereas a locally uniform finite limit of entire functions on every
compact subset of $\C$ is entire.  The operator convergence proved here does not require this
problematic entire-function convergence.  It works directly with resolvents and yields the
spectral conclusion that the proposal was intended to motivate.

Strong resolvent convergence is weaker than locally uniform convergence of normalized
characteristic functions.  This weakness is appropriate: it proves that every target xi ordinate is
seen by the finite spectra but does not force a single entire normalization, and it does not exclude
additional finite eigenvalues.  A stronger meromorphic Liv\v{s}ic-function convergence would
require source estimates not supplied by the abstract phase-locking theorem.

\section{Significance, novelty boundary, and open problems}

The contribution has two layers.

The universal layer is the nested-exhaustion phase-locking law.  Stable graph exhaustion
automatically converges the deficiency spaces; polar transport then constructs and classifies
self-adjoint extension limits.  Bjerg's fixed-space theorem explains how convergent boundary maps
control convergent extensions \cite{Bjerg2024}.  Our theorem adds an automatic and canonical choice
of those maps in varying nested spaces.

The arithmetic layer resolves Suzuki's stated finite-aperture spectral-limit problem under RH.
The missing phase exists for structural reasons, and its only asymptotic degree of freedom is one
unit complex number.  Every xi ordinate is approximated by real finite-aperture characteristic
zeros.

Several frontiers remain.

\begin{enumerate}
 \item The theorem does not prove RH.  Its Suzuki specialization begins with the positive global
 Weil Hilbert space and the self-adjoint target generator, both available under RH.
 \item The polar formula is canonical but not an explicit prime-by-prime scalar formula for
 $\theta(a)$.  Deriving such a formula would connect boundary locking to Suzuki's finite arithmetic
 kernel.
 \item Generalized strong resolvent convergence allows spectral pollution.  Norm-resolvent or
 meromorphic Liv\v{s}ic convergence would be needed for two-sided spectral convergence.
 \item The theorem selects the target extension once it is specified.  An unconditional source
 principle selecting the global positive Weil space would go beyond this paper and would amount to
 progress on RH itself.
\end{enumerate}

The operator theorem is exact and unconditional.  Priority for its particular nested-space
formulation and the Suzuki application requires independent specialist review.

\begin{thebibliography}{99}

\bibitem{Bjerg2024}
A.~Bjerg,
\emph{Convergence of operators with deficiency indices $(k,k)$ and of their self-adjoint
extensions},
Ann. Henri Poincar\'e \textbf{25} (2024);
2995--3007; doi:10.1007/s00023-023-01397-9.

\bibitem{ReedSimonI}
M.~Reed and B.~Simon,
\emph{Methods of Modern Mathematical Physics. I: Functional Analysis},
second edition, Academic Press, 1980.

\bibitem{ReedSimonII}
M.~Reed and B.~Simon,
\emph{Methods of Modern Mathematical Physics. II: Fourier Analysis, Self-Adjointness},
Academic Press, 1975.

\bibitem{Suzuki2023}
M.~Suzuki,
\emph{Aspects of the screw function corresponding to the Riemann zeta-function},
J. London Math. Soc. \textbf{108} (2023), 1448--1487.

\bibitem{Suzuki2026}
M.~Suzuki,
\emph{Weil's quadratic form via the screw function},
arXiv:2606.09096, 2026.

\bibitem{TeschlEtAl2026}
G.~Teschl, Y.~Wang, B.~Xie, and Z.~Zhou,
\emph{On generalized strong and norm resolvent convergence},
arXiv:2601.10476, 2026.

\end{thebibliography}

\end{document}
